\exercisesection{}

\begin{enumerate}
\def\labelenumi{\arabic{enumi})}
\item
  设 G 为局部紧群。设 \(\pi\) 为 G 在复希尔伯特空间 E 中的酉表示，且 \(x \in \mathrm{E}\) 。我们有 \(\pi(g)x = x\) （对所有 \(g \in \mathrm{G}\) ），当且仅当正定函数 \(g \mapsto \langle x | \pi(g)x \rangle\) 为常数。
\item
  设 G 为局部紧群。对 G 上的每个正定函数 \(\varphi\) ，我们用 \((\pi_{\varphi}, x_{\varphi})\) 表示 \(\varphi\) 的一个循环希尔伯特表示。
\end{enumerate}

\begin{enumerate}
\def\labelenumi{\alph{enumi})}
\item
  设 \(\varphi_1, \varphi_2\) 与 \(\varphi_3\) 为 G 上的正定函数。若 \(\varphi_3 = \varphi_1 + \varphi_2\) ，则存在从 \(\pi_{\varphi_1}\) 到 \(\pi_{\varphi_3}\) 的单射 G-态射。
\item
  若 \(\varphi\) 是非零的常值正定函数，则 \(\pi_{\varphi}\) 是一维的平凡表示。
\item
  对 G 在希尔伯特空间 E 中的每个酉表示 \(\pi\) 以及每个 \(x \in \mathrm{E}\) ，存在从 \(\pi_{\varphi}\) 到 \(\pi\) 的单射 G-态射，其中 \(\varphi(g) = \langle x \mid \pi(g)x \rangle\) 。
\item
  设 H 为 G 的闭子群，\(\varphi\) 为 G 上的正定函数。用 \(\psi\) 表示 \(\varphi\) 在 H 上的限制；它是 H 上的正定函数，且存在从 \(\pi_{\psi}\) 到 \(\operatorname{Res}_{\mathrm{H}}^{\mathrm{H}}(\pi)\) 的单射 H-态射。
\end{enumerate}

\begin{enumerate}
\def\labelenumi{\arabic{enumi})}
\setcounter{enumi}{2}
\tightlist
\item
  \begin{enumerate}
  \def\labelenumii{\alph{enumii})}
  \tightlist
  \item
  给出一个对合巴拿赫代数的例子，它具有一个不连续的正线性泛函。
  \end{enumerate}
\end{enumerate}

\begin{enumerate}
\def\labelenumi{\alph{enumi})}
\setcounter{enumi}{1}
\tightlist
\item
  给出一个对合巴拿赫代数 A、A 上一个连续的正线性泛函 \(\lambda\) ，以及希尔伯特实现 \((\mathrm{E}_{1}, e_{1}, \varphi_{1})\) 与 \((\mathrm{E}_{2}, e_{2}, \varphi_{2})\) （均为 \(\lambda\) 的实现）的例子，使得 \((\mathrm{E}_{1}, e_{1}, \varphi_{1})\) 是循环希尔伯特实现，并且 V, p.~440 命题 6 中唯一的线性映射 u 不满足 \(u(e_{1}) = e_{2}\) 。
\end{enumerate}

\begin{enumerate}
\def\labelenumi{\arabic{enumi})}
\setcounter{enumi}{3}
\tightlist
\item
  设 A 与 B 为星代数。从 A 到 B 的线性映射 u 称为正的，如果 \(u(\mathrm{A}_{+}) \subset \mathrm{B}_{+}\) ；称为完全正的，如果对每个整数 \(k \in N\) ，映射 \(u_{k} = 1_{C^{k}} \otimes u: C^{k} \otimes A \to C^{k} \otimes B\) 是正的。
\end{enumerate}

\begin{enumerate}
\def\labelenumi{\alph{enumi})}
\tightlist
\item
  设 A 为含幺星代数，E 与 F 为复希尔伯特空间，v 为从 E 到 F 的线性映射，\(\varphi\) 为从 A 到 \(\mathcal{L}(\mathrm{F})\) 的对合代数的保幺态射，使得 \(\|\varphi\|^{2} = \|\varphi(1)\|\) 。映射 \(u: A \to \mathcal{L}(E)\) （由 \(u(a) = v^{*} \circ \varphi(a) \circ v\) 定义）是完全正的。
\end{enumerate}

本习题的目标是证明这一断言的逆命题（“斯廷斯普龙定理”）。设 A 为含幺星代数，E 为复希尔伯特空间，\(u \colon \mathrm{A} \to \mathcal{L}(\mathrm{E})\) 为完全正线性映射。

\begin{enumerate}
\def\labelenumi{\alph{enumi})}
\setcounter{enumi}{1}
\tightlist
\item
  设 \(F = A \otimes E\) 。F 上存在一个正的半双线性型 q，使得
\end{enumerate}

\[
q (a \otimes x, b \otimes y) = \langle x | \varphi (a ^ {*} b) y \rangle
\]

对所有 \((a,b)\in \mathrm{A}^2\) 与 \((x,y)\in \mathrm{E}^2\) 成立。

设 \(\widehat{\mathbf{F}}\) 是 \(\mathbf{F}\) 关于半双线性型 \(q\) 的分离完备希尔伯特空间。

\begin{enumerate}
\def\labelenumi{\alph{enumi})}
\setcounter{enumi}{2}
\tightlist
\item
  设 \(b \in A\) 。F 上使 \(a \otimes x \mapsto ba \otimes x\) （对所有 \(a \in A\) 与 \(x \in E\) ）成立的自同态，通过过渡到商诱导出 \(\varphi(b)\) （\(\widehat{F}\) 的一个自同态），满足 \(\|\varphi(b)\|\leqslant\|b\|\) ；映射 \(\varphi\colon A\to\mathcal{L}(F)\) 是对合代数的保幺态射。
\end{enumerate}

\begin{enumerate}
\def\labelenumi{\alph{enumi})}
\setcounter{enumi}{3}
\tightlist
\item
  设 \(v\) 为从 E 到 \(\widehat{\mathrm{F}}\) 的线性映射，使得 \(v(x)\) 是 \(1 \otimes x\) 所在的类（对所有 \(x \in \mathrm{E}\) ）。我们有 \(\|v\|^2 = \| \varphi(1)\|\) ，且 \(u(a) = v^* \circ \varphi(a) \circ v\) 对所有 \(a \in \mathrm{A}\) 成立。
\end{enumerate}

\begin{enumerate}
\def\labelenumi{\arabic{enumi})}
\setcounter{enumi}{4}
\tightlist
\item
  设 X 为分离拓扑空间。连续函数 \(f: X \times X \to R\) 称为条件负定型的，如果它满足下列条件
\end{enumerate}

\begin{enumerate}
\def\labelenumi{(\roman{enumi})}
\item
  对每个 \(x \in X\) ，有 \(f(x, x) = 0\) ；
\item
  对所有 \(x\) 与 \(y\) （属于 \(X\) ），我们有 \(f(x,y) = f(y,x)\) ；
\item
  对每个整数 \(n \in N\) 、X 中元素的每个族 \((x_{i})_{0 \leqslant i \leqslant n}\) 以及实数的每个族 \((t_{i})_{0 \leqslant i \leqslant n}\) （满足
\end{enumerate}

\[
\sum_ {i = 1} ^ {n} t _ {i} = 0,
\]

），我们有

\[
\sum_ {i = 0} ^ {n} \sum_ {j = 0} ^ {n} t _ {i} t _ {j} f (x _ {i}, x _ {j}) \leqslant 0.
\]

我们用 Noy \(_{-}\) (X) 表示 X × X 上条件负定型函数构成的集合。

\begin{enumerate}
\def\labelenumi{\alph{enumi})}
\item
  设 E 为实希尔伯特空间，\(\varphi\colon\mathrm{X}\to\mathrm{E}\) 为连续映射。函数 \(f(x,y)=\|\varphi(x)-\varphi(y)\|^{2}\) 是条件负定型的。
\item
  设 \(f\in\mathrm{Noy}_{-}(\mathrm{X})\) 且 \(x_{0}\in\mathrm{X}\) 。存在一个实希尔伯特空间 E 与一个连续映射 \(\varphi\colon\mathrm{X}\to\mathrm{E}\) ，使得
  \begin{enumerate}
  \def\labelenumii{\alph{enumii})}
  \tightlist
  \item
    对每个 \((x,y)\in \mathrm{X}\times \mathrm{X}\) ，有 \(f(x,y) = \| \varphi (x) - \varphi (y)\| ^2\) ；
  \item
    族 \((\varphi(x)-\varphi(x_{0}))_{x\in\mathrm{X}}\) 在 E 中是完全的。
  \end{enumerate}
\end{enumerate}

我们称 \((\mathrm{E},\varphi)\) 为 f 的一个循环希尔伯特表示。

\begin{enumerate}
\def\labelenumi{\alph{enumi})}
\setcounter{enumi}{2}
\item
  设 \(f \in \mathrm{Noy}_{-}(\mathrm{X})\) 且 \(x_0 \in \mathrm{X}\) 。设 \((\mathrm{E}, \varphi)\) 与 \((\mathrm{E}', \varphi')\) 为 \(f\) 的循环希尔伯特表示。存在唯一的仿射等距 \(u: \mathrm{E} \to \mathrm{E}'\) ，使得 \(\varphi' = u \circ \varphi\) 。
\item
  集合 \(\mathrm{Noy}_{-}(\mathrm{X})\) 是 \(\mathcal{C}(\mathrm{X}\times \mathrm{X};\mathbf{R})\) 中的一个凸锥。对每个普遍正核 \(g\in \mathrm{Noy}_{+}(\mathrm{X})\) （\(g\) 取实值），映射 \(f(x,y) = g(x,x) - g(x,y)\) 属于 \(\mathrm{Noy}_{-}(\mathrm{X})\) 。
\item
  设 \(\mathfrak{F}\) 为 \(\mathrm{Noy}_{-}(\mathrm{X})\) 上逐点收敛于 \(f\colon \mathrm{X}\times \mathrm{X}\to \mathbf{R}\) 的滤子。若 \(f\) 连续，则 \(f\in \mathrm{Noy}_{-}(\mathrm{X})\) 。
\item
  设 \(f: X \times X \to R\) 为满足上述条件 (i) 与 (ii) 的连续函数。设 \(x_{0} \in X\)。设 \(g: X \times X \to R\) 满足 \(g(x, y) = f(x, x_{0}) + f(y, x_{0}) - f(x, y)\) 。则 \(f \in \text{Noy}_{-}(X)\) 当且仅当 \(g \in \text{Noy}_{+}(X)\) 。
\item
  假设 X 非空。设 \(f: \mathrm{X} \times \mathrm{X} \to \mathbf{R}\) 为满足上述条件 (i) 与 (ii) 的连续函数。则 \(f \in \mathrm{Noy}_{-}(\mathrm{X})\) 当且仅当 \(\exp(-tf) \in \mathrm{Noy}_{+}(\mathrm{X})\) 对所有 \(t \in \mathbf{R}_{+}\) 成立（“舍恩贝格定理”）。（为证 \(\exp(-f)\) 在 \(f \in \mathrm{Noy}_{-}(\mathrm{X})\) 时是正定型的，固定 \(x_0 \in \mathrm{X}\) ；证明
\end{enumerate}

\[
g (x, y) = \exp (f (x, x _ {0}) + f (y, x _ {0}) - f (x, y))
\]

以及

\[
h (x, y) = \exp (- f (x, x _ {0})) \exp (- f (y, x _ {0}))
\]

都属于 Noy \(_{+}\) (X)。）

若 G 是拓扑群，我们称函数 \(\varphi \in \mathcal{C}(\mathrm{G})\) 是负定型的，如果函数 \(f\) 是 \(\mathrm{G} \times \mathrm{G}\) 上由 \(f(g,h) = \varphi(h^{-1}g)\) 定义的条件负定核。

\begin{enumerate}
\def\labelenumi{\arabic{enumi})}
\setcounter{enumi}{5}
\tightlist
\item
  设 G 为离散拓扑群。设 A 为 \(\mathcal{L}(\ell^{2}(\mathrm{G}))\) 中由左正则表示 \(\lambda\) 的像生成的对合子代数。这是一个含幺星代数。
\end{enumerate}

对每个正定型函数 \(\varphi\in\mathrm{Pos}(\mathrm{G})\) ，存在唯一的从 A 到 A 的完全正线性映射 \(u_{\varphi}\) （参见习题 4），使得 \(u_{\varphi}(\boldsymbol{\lambda}(g))=\varphi(g)\boldsymbol{\lambda}(g)\) 对所有 \(g\in\mathrm{G}\) 成立。（利用 \(\varphi\) 的一个循环希尔伯特表示。）

\begin{enumerate}
\def\labelenumi{\arabic{enumi})}
\setcounter{enumi}{6}
\tightlist
\item
  设 I 为有限集，\(G = F(I)\) 为在 I 上构造的自由群（A, I, p.~84）。设 A 为 \(\mathcal{L}(\ell^{2}(G))\) 中由左正则表示 \(\lambda\) 的像生成的对合子代数。
\end{enumerate}

\begin{enumerate}
\def\labelenumi{\alph{enumi})}
\item
  对每个 \(g \in G\) ，用 \(|g|\) 表示 g 的唯一既约分解的长度（A, I, p.~146, 习题 19）。映射 \(d(g, h) = |h^{-1}g|\) 是 G 上的一个距离。
\item
  设 \(\lambda \in \mathbf{R}_{+}^{*}\) 。G 上由 \(\psi_{\lambda}(g) = \exp(-\lambda|g|)\) 定义的函数是正定函数。（证明函数 \(g \mapsto |g|\) 是负定型函数，参见习题~5。）
\item
  对每个整数 \(k \in N\) ，设 \(G_{k}\) 为由 \(g \in G\) （满足 \(|g| = k\) ）组成的集合，并设 \(\varphi_{k}\) 为 \(G_{k}\) 的特征函数。设 f 与 g 为 G 上的函数，其支集分别含于 \(G_{k}\) 与 \(G_{l}\) 。对每个 \(m \in N\) ，我们有 \((f * g)\varphi_{m} = 0\) ，除非 \(k + l - m\) 为偶数且 \(|k - l| \leqslant m \leqslant k + l\) ；在这一情形，我们有
\end{enumerate}

\[
\left\| (f * g) \varphi_ {m} \right\| \leqslant \| f \| \| g \|.
\]

\begin{enumerate}
\def\labelenumi{\alph{enumi})}
\setcounter{enumi}{3}
\tightlist
\item
  设 f 为 G 上具有有限支集的函数。我们有
\end{enumerate}

\[
\| \boldsymbol {\lambda} (f) \| \leqslant \sum_ {n \in \mathbf {N}} (n + 1) \| f \varphi_ {n} \|,
\]

以及

\[
\left\| \boldsymbol {\lambda} (f) \right\| \leqslant 2 \left(\sum_ {g \in G} (1 + | g |) ^ {4} | f (g) | ^ {2}\right) ^ {1 / 2}.
\]

\begin{enumerate}
\def\labelenumi{\alph{enumi})}
\setcounter{enumi}{4}
\item
  设 \(\psi\) 为 G 上的函数，使得 \(M = \sup_{g \in G} |\psi(g)| (1 + |g|)^2\) 有限。存在唯一的从 A 到 A 的线性映射 \(u_\psi\) ，使得 \(u_\psi(\boldsymbol{\lambda}(g)) = \psi(g)\boldsymbol{\lambda}(g)\) 对所有 \(g \in G\) 成立。我们有 \(\|u_\psi\| \leqslant 2M\) 。
\item
  设 X 为 G 上由带有限支集的函数 \(\psi\) 组成的集合，其中 \(\|u_{\psi}\|\leqslant1\) 。存在 X 上的一个滤子 \(\mathfrak{F}\) ，使得
\end{enumerate}

\[
\lim _ {\psi , \mathfrak {F}} u _ {\psi} = 1 _ {\mathrm{A}}
\]

在赋予逐点收敛拓扑的空间 \(\mathcal{L}(\mathrm{A})\) 中。特别地，A 的底巴拿赫空间具有逼近性质。（考虑函数 \(\psi_{\lambda,n}(g) = (\varphi_0 + \dots + \varphi_n)\psi_\lambda\) ，其中 \(\lambda \in \mathbf{R}_+^*\) ，\(n \in \mathbf{N}\) ，并利用上一习题。）

\begin{enumerate}
\def\labelenumi{\arabic{enumi})}
\setcounter{enumi}{7}
\tightlist
\item
  我们把群 \(\mathbf{R}\) 与它的对偶群通过映射 \((x,y)\mapsto\) \(\exp (i\pi xy)\) 等同起来。设 \(\alpha \in \mathbf{R}_{+}\) 。
\end{enumerate}

\begin{enumerate}
\def\labelenumi{\alph{enumi})}
\item
  \(\mathbf{R}\) 上存在其傅里叶变换为函数 \(f_{\alpha}(x) = \exp(-|x|^{\alpha})\) 的正测度，当且仅当 \(\alpha \in ]0,2]\) 。在这一情形，我们把这一测度记作 \(\mu_{\alpha}\) ，它是唯一的；称之为参数 \(\alpha\) 的稳定测度。
\item
  测度 \(\mu_{\alpha}\) 的总质量为 1；R 上的恒等函数 f 关于 \(\mu_{\alpha}\) 可积，当且仅当 \(\alpha = 2\) 。
\end{enumerate}

\begin{enumerate}
\def\labelenumi{\arabic{enumi})}
\setcounter{enumi}{8}
\tightlist
\item
  本习题的目标是给出 V, p.~455 定理 5 的另一个证明。设 G 为局部紧交换拓扑群，带有哈尔测度 \(\mu\) 。我们用 \(\widehat{f}\) 表示 \(f \in \mathrm{L}^1(\mathrm{G})\) 的傅里叶变换，并用 \(\mathcal{F}(\mathrm{G}) \subset \mathcal{C}_0(\widehat{\mathrm{G}})\) 表示 \(\mathrm{L}^1(\mathrm{G})\) 上傅里叶变换的像。
\end{enumerate}

\begin{enumerate}
\def\labelenumi{\alph{enumi})}
\item
  空间 \(\mathcal{F}(\mathrm{G})\) 是 \(\mathcal{C}_0(\mathrm{G})\) 的稠密子空间。
\item
  对每个 \(f \in \mathrm{L}^{1}(\mathrm{G})\) ，f 在巴拿赫代数 \(\mathrm{L}^{1}(\mathrm{G})\) 中的谱半径等于 \(\|\widehat{f}\|_{\infty}\) 。
\item
  设 \(\varphi\in\mathrm{Pos}(\mathrm{G})\) 。它是 \(\mathrm{G}\) 上一致连续且有界的函数；对每个 \(f\in\mathrm{L}^{1}(\mathrm{G})\) ，我们有
\end{enumerate}

\[
\int_ {G} \int_ {G} f (x) \overline {{f (y)}} \varphi (x - y) d (\mu \otimes \mu) (x, y) \geqslant 0.
\]

\begin{enumerate}
\def\labelenumi{\alph{enumi})}
\setcounter{enumi}{3}
\tightlist
\item
  设 \(\lambda\) 为 \(\mathrm{L}^{1}(\mathrm{G})\) 上由 \(\varphi\) 确定的连续线性型，并令 \(b(f,g)=\lambda(f^{*}*g)\) （对所有 \((f,g)\in\mathrm{L}^{1}(\mathrm{G})\times\mathrm{L}^{1}(\mathrm{G})\) ）。映射 \(b\) 是 \(\mathrm{L}^{1}(\mathrm{G})\) 上的一个正埃尔米特型。
\end{enumerate}

\begin{enumerate}
\def\labelenumi{\alph{enumi})}
\setcounter{enumi}{4}
\tightlist
\item
  对每个 \(f \in \mathrm{L}^1(\mathrm{G})\) ，我们有 \(|\lambda(f)|^2 \leqslant b(f, f)\) 。
\end{enumerate}

\begin{enumerate}
\def\labelenumi{\alph{enumi})}
\setcounter{enumi}{5}
\tightlist
\item
  设 \(f \in \mathrm{L}^{1}(\mathrm{G})\) ，\(g = f^{*} * f \in \mathcal{C}_{b}(\mathrm{G})\) 。定义 \(g_{1} = g\) 与 \(g_{n+1} = g_{n} * g\) （对每个整数 \(n \geqslant 1\) ）。对每个整数 \(n \geqslant 1\) ，我们有 \(|\lambda(f)| \leqslant \|g_{n}\|_{1}^{2^{-n}}\) 。
\end{enumerate}

\begin{enumerate}
\def\labelenumi{\alph{enumi})}
\setcounter{enumi}{6}
\tightlist
\item
  对每个 \(f \in \mathrm{L}^{1}(\mathrm{G})\) ，我们有 \(|\lambda(f)| \leqslant \| \widehat{f} \|_{\infty}\) 。存在连续线性映射 \(\widehat{\lambda}\) （定义在 \(\mathcal{F}(\mathrm{G})\) 上），使得 \(\lambda(f) = \widehat{\lambda}(\widehat{f})\) 对所有 \(f \in \mathrm{L}^{1}(\mathrm{G})\) 成立。
\end{enumerate}

\begin{enumerate}
\def\labelenumi{\alph{enumi})}
\setcounter{enumi}{7}
\tightlist
\item
  存在有界测度 \(\nu\) （在 \(\widehat{G}\) 上），使得 \(\lambda(f) = \nu(\widehat{f})\) 对每个 \(f \in \mathrm{L}^{1}(\mathrm{G})\) 成立；测度 \(\nu\) 是正的，并且有
\end{enumerate}

\[
\varphi (x) = \int_ {\widehat {\mathrm{G}}} \langle \chi , x \rangle d \nu (\chi)
\]

对所有 \(x \in G\) 成立。

\begin{enumerate}
\def\labelenumi{\arabic{enumi})}
\setcounter{enumi}{9}
\tightlist
\item
  \begin{enumerate}
  \def\labelenumii{\alph{enumii})}
  \tightlist
  \item
    设 \(\mathrm{G} = \mathbf{R}\) 。集合 \(\operatorname{Pos}_1(\mathrm{G})\) 在弱拓扑下不是紧的。
  \end{enumerate}
\end{enumerate}

\begin{enumerate}
\def\labelenumi{\alph{enumi})}
\setcounter{enumi}{1}
\tightlist
\item
  设 G = R/Z。\(\operatorname{Pos}_{\leqslant1}(\mathbf{G})\) 上的弱拓扑与紧收敛拓扑不一致。
\end{enumerate}

\begin{enumerate}
\def\labelenumi{\arabic{enumi})}
\setcounter{enumi}{10}
\tightlist
\item
  设 G 为局部紧群。
\end{enumerate}

设 \(\varrho\) 与 \(\pi\) 为 G 的酉表示。我们称 \(\varrho\) 弱包含于 \(\pi\) ，并记作 \(\varrho \leqslant \pi\) ，如果 \(\varrho\) 的对角矩阵系数所成的集合包含于 \(\pi\) 的对角矩阵系数的有限和所成的集合在赋予紧收敛拓扑的空间 \(\mathcal{C}(\mathrm{G})\) 中的闭包内。

\begin{enumerate}
\def\labelenumi{\alph{enumi})}
\item
  关系“\(\varrho\) 与 \(\pi\) 是 G 的酉表示且 \(\varrho \leqslant \pi\)”是一个预序关系。若 \(\varrho'\) （相应地 \(\pi'\) ）同构于 \(\varrho\) （相应地 \(\pi\) ），则 \(\varrho \leqslant \pi\) 当且仅当 \(\varrho' \leqslant \pi'\) 。
\item
  若 \(\varrho\) 同构于 \(\pi\) 的一个子表示，则 \(\varrho \leqslant \pi\) 。若 I 与 J 是非空有限集，且若 \(\varrho \leqslant \pi\) ，则 \(\bigoplus_{i \in I} \varrho \leqslant \bigoplus_{j \in J} \pi\) 。
\item
  设 X 为 \(\varrho\) 的空间 E 中生成 \(\pi\) 的一个子集。我们有 \(\varrho \leqslant \pi\) ，当且仅当每个矩阵系数 \(g \mapsto \langle x | \varrho(g)x \rangle\) （带 \(x \in X\) ）都属于 \(\pi\) 的对角矩阵系数的有限和所成的集合在赋予紧收敛拓扑的空间 \(\mathcal{C}(\mathrm{G})\) 中的闭包。（证明 E 中满足这一性质的元素所成的集合 Y 是 E 的一个闭向量子空间。）
\item
  设 H 为 G 的闭子群。设 \(\pi\) 与 \(\varrho\) 为 G 的酉表示，使得 \(\varrho \leqslant \pi\) ；则我们有 \(\operatorname{Res}_{\mathrm{H}}^{\mathrm{G}}(\varrho) \leqslant \operatorname{Res}_{\mathrm{H}}^{\mathrm{G}}(\pi)\) 。
\item
  设 H 为 G 的闭子群，\(\nu\) 为 H 上的哈尔测度。设 \(\pi\) 与 \(\varrho\) 为 H 的酉表示，使得 \(\varrho \leqslant \pi\) ；则我们有 \(\operatorname{Ind}_{\mathrm{H}}^{\mathrm{G}}(\varrho) \leqslant \operatorname{Ind}_{\mathrm{H}}^{\mathrm{G}}(\pi)\) 。（令 \(\varrho' = \operatorname{Ind}_{\mathrm{H}}^{\mathrm{G}}(\varrho)\) ，\(\pi' = \operatorname{Ind}_{\mathrm{H}}^{\mathrm{G}}(\pi)\) ；利用 (c)，化归为逼近形如 \(g \mapsto \langle f | \varrho'(g)f \rangle\) 的矩阵系数，其中 f 形如
\end{enumerate}

\[
f (x) = \int_ {\mathrm{H}} f (x h) \varrho (h) y d \nu (y)
\]

（其中 \(\varphi\in\mathcal{K}(\mathrm{G})\) ，\(y\) 属于 \(\varrho\) 的空间。）

\begin{enumerate}
\def\labelenumi{\arabic{enumi})}
\setcounter{enumi}{11}
\item
  设 G = R，\(\pi\) 为 G 在 \(\mathrm{L}^{2}(\mathbf{R})\) 中的正则表示。证明 \(\chi \leqslant \pi\) 对每个特征标 \(\chi \in \widehat{G}\) 成立。
\item
  设 \(\varrho\) 为 G 在希尔伯特空间 E 中的不可约酉表示，\(\pi\) 为 G 在 F 中的酉表示。设 \(x \in E\) 为范数 1 的向量。假设 \(\varrho \leqslant \pi\) （参见习题~11）。
\end{enumerate}

我们用 \(A \subset \mathcal{C}(G)\) 表示 \(\pi\) 的规范化的对角矩阵系数所成的集合，用 \(\widetilde{A}\) 表示 A 在 \(L^{\infty}(G)\) 中关于弱拓扑的闭凸包。

设 \(x \in E\) 为范数 1 的向量，\(f: g \mapsto \langle x \mid \varrho(g)x \rangle\) 为相伴的对角矩阵系数。

\begin{enumerate}
\def\labelenumi{\alph{enumi})}
\item
  集合 \(\widetilde{\mathrm{A}}\) 是 \(\mathrm{L}^{\infty}(\mathrm{G})\) 中的一个紧凸集，而 \(f\) 是 \(\widetilde{\mathrm{A}}\) 的一个极点。
\item
  函数 f 属于 A 在 \(\widetilde{A}\) 中关于弱拓扑的闭包。
\item
  函数 f 属于 A 在赋予紧收敛拓扑的 \(\mathcal{C}(G)\) 中的闭包。
\item
  G 在 C 中的平凡表示弱包含于 \(\pi\) ，当且仅当对 G 的每个紧子集 C 与每个实数 \(\varepsilon > 0\) ，存在范数为 1 的 \(x \in F\) ，使得
\end{enumerate}

\[
\sup _ {g \in \mathrm{C}} \| \pi (g) x - x \| <   \varepsilon .
\]

\begin{enumerate}
\def\labelenumi{\arabic{enumi})}
\setcounter{enumi}{13}
\tightlist
\item
  设 G 为局部紧群。我们称 G 具有卡日丹性质 ]T[ \(^{(1)}\) ，如果 G 的酉表示 \(\pi\) 具有非零的平凡子表示 \(^{(2)}\) ，当且仅当 G 在 C 上的平凡表示 \(\mathbf{1}_{\text{G}}\) 弱包含于 \(\pi\) 。
\end{enumerate}

\begin{enumerate}
\def\labelenumi{\alph{enumi})}
\item
  设 G 为具有性质 ]T[ 的局部紧群。设 H 为 G 中由 G 的紧子集生成的开子群所成的集合。群 G 是诸子群 H ∈ H 的并。
\item
  设 \(\pi\) 为 G 的这样一个表示，它是 G 在 G/H（\(H \in H\) ）上的拟正则表示的希尔伯特和（参见 V, p. 486 的习题 2）。我们有 \(1_{G} \leqslant \pi\) 。
\item
  群 G 由 G 的一个紧子集生成。特别地，若 G 是离散的，则 G 由 G 的一个有限子集生成。
\end{enumerate}

\begin{enumerate}
\def\labelenumi{\arabic{enumi})}
\setcounter{enumi}{14}
\tightlist
\item
  设 G 为局部紧群，C 为 G 的子集，\(\varepsilon > 0\) 。我们称 \((\mathrm{C}, \varepsilon)\) 是 G 的一个卡日丹对，如果对 G 在希尔伯特空间 E 中的每个酉表示 \(\pi\) ，存在向量 \(x \in E\)
\end{enumerate}

使得

\[
\sup _ {g \in \mathrm{C}} \| \pi (g) x - x \| <   \varepsilon \| x \|,
\]

  便蕴含 \(\pi\) 有一个非零的平凡子表示。我们称 C 为卡日丹集，如果存在 \(\varepsilon > 0\) 使得 \((C, \varepsilon)\) 是一个卡日丹对。

\begin{enumerate}
\def\labelenumi{\alph{enumi})}
\tightlist
\item
  设 \((\mathrm{C}, \varepsilon)\) 为 G 的一个卡日丹对。设 \(\pi\) 为 G 在希尔伯特空间 E 中的酉表示。设 \(x \in E\) 非零且 \(\delta > 0\) ，使得
\end{enumerate}

\[
\sup _ {g \in \mathrm{C}} \| \pi (g) x - x \| <   \delta \varepsilon \| x \|.
\]

则我们有 \(\|p(x)-x\|\leqslant\delta\|x\|\) ，其中 p 是 E 的正交投影算子，其像为 G 不变向量所成的空间。

\begin{enumerate}
\def\labelenumi{\alph{enumi})}
\setcounter{enumi}{1}
\item
  局部紧群 G 具有性质 ]T[，当且仅当存在一个卡日丹对 (C,ε)。（为证性质 ]T[ 蕴含这一断言，用反证法，考虑诸表示 π \(_{C,\varepsilon}\) 的希尔伯特和表示，这些表示由 G 的紧子集 C 与 ε > 0 参数化，使得 π \(_{C,\varepsilon}\) 不含非平凡子表示，但存在范数为 1 的 x \(_{C,\varepsilon}\) 满足 sup \(_{g\in\mathbf{C}}\|\pi(g)x_{\mathrm{C},\varepsilon}-x_{\mathrm{C},\varepsilon}\|<\varepsilon.\)
\item
  设 G 为具有性质 ]T[ 的局部紧群。G 的每个生成 G 的紧子集 C 都是卡日丹集。
\item
  设 G 为具有性质 ]T[ 的局部紧群，C 为 G 的一个紧子集，其内部非空且是卡日丹集。则 C 生成 G。特别地，若 G 是离散的，则卡日丹集恰好是 G 的有限生成子集。（由 C 生成的 G 的子群 H 是开的；考虑 G 在 G/H 上的拟正则表示 \(\pi\) ，并证明 H 的特征函数 \(\varphi\) （H \(\in\) G/H）在 \(\pi\) 的不变向量空间上的投影等于 \(\varphi\) 。）
\end{enumerate}

\begin{enumerate}
\def\labelenumi{\arabic{enumi})}
\setcounter{enumi}{15}
\item
  设 G 为在无穷远处可数的局部紧群，H 为 G 的闭子群。若 G 具有性质 ]T[，且若 G/H 容许一个有界的 G 不变测度，则 H 具有性质 ]T[。（利用 V, p.~492 的习题 20。）
\item
  设 \(n \geqslant 2\) 为整数。我们考虑实李群 \(\mathbf{G} = \mathbf{R}^n \rtimes \mathbf{SL}(n, \mathbf{R})\) ，并用 H 表示 \(\mathbf{R}^n \times \{e\}\) （\(\mathbf{G}\) 的子群），它与 \(\mathbf{R}^n\) 等同。
\end{enumerate}

设 \(\pi\) 为 G 在复希尔伯特空间 E 中的酉表示。假设对 G 的每个紧子集 C 与每个 \(\varepsilon > 0\) ，存在非零的 \(x \in E\) ，使得

\[
\sup _ {g \in \mathrm{C}} \| \pi (g) x - x \| <   \varepsilon \| x \|.
\]

\begin{enumerate}
\def\labelenumi{\alph{enumi})}
\item
  存在 E 中由范数为 1 的向量构成的序列 \((x_{m})\) ，使得函数 \(\varphi_{m}(g)=\langle x_{m}|\pi(g)x_{m}\rangle\) 在赋予紧收敛拓扑的 \(\mathcal{C}(\mathrm{G})\) 中收敛于常值函数 1。（考虑一个其并等于 G 的紧子集序列 \((\mathrm{C}_{m})_{m\in\mathbf{N}}\) ，并取 \(x_{m}\) 使其对 \((\mathrm{C}_{m},1/(m+1))\) 满足上述条件。）
\item
  对每个 \(m \in N\) ，存在总质量为 1 的正测度 \(\mu_{m}\) （在 \(R^{n}\) 上），其傅里叶变换与 \(\varphi_{m}\) 在 H 上的限制等同。
\item
  假设 \(\mu_{m}(\{0\}) = 0\) 对所有 \(m \in N\) 成立。这时把 \(\mu_{m}\) 与 \(R^{n} - \{0\}\) 上总质量为 1 的正测度等同起来，并用 \(\nu_{m}\) 表示 \(\mu_{m}\) 在射影空间 \(\mathbf{P}_{n-1}(\mathbf{R})\) 上的像测度。存在子序列 \((\nu_{n_{k}})_{k \in \mathbf{N}}\) ，它狭收敛于总质量为 1 的正测度 \(\nu\) （在 \(\mathbf{P}_{n-1}(\mathbf{R})\) 上）；测度 \(\nu\) 是 \(\mathbf{SL}(n, \mathbf{R})\) -不变的。由此导出矛盾。
\item
  设 \(m \in N\) 满足 \(\mu_{m}(\{0\}) > 0\) 。设 \(\psi\) 为 \(\varphi_{m}\) 在 H 上的限制。存在 c > 0，使得 \(\psi = c + \widetilde{\psi}\) ，其中 \(\widetilde{\psi}\) 是 \(R^{n}\) 上的正定型函数。
\item
  表示 \(\pi\) 含有 H 不变的向量。（利用习题 2。）我们称 (G,H) 具有相对性质 ]T[。
\end{enumerate}

\begin{enumerate}
\def\labelenumi{\arabic{enumi})}
\setcounter{enumi}{17}
\tightlist
\item
  设 \(\mathrm{G} = \mathbf{SL}(2, \mathbf{R})\) 。用 N 与 P 表示子群
\end{enumerate}

\[
\mathrm{N} = \Big \{\left( \begin{array}{c c} 1 & t \\ 0 & 1 \end{array} \right) | t \in \mathbf {R} \Big \}, \qquad \mathrm{P} = \Big \{\left( \begin{array}{c c} a & t \\ 0 & a ^ {- 1} \end{array} \right) | a \in \mathbf {R} _ {*}, t \in \mathbf {R} \Big \}
\]

（G 的子群）。设 \(\pi\) 为 G 在复希尔伯特空间 E 中的酉表示。假设存在 \(x \in \mathrm{E}\) ，使得 \(\pi(n)x = x\) 对所有 \(n \in \mathrm{N}\) 成立。

\begin{enumerate}
\def\labelenumi{\alph{enumi})}
\tightlist
\item
  齐性空间 \(\mathbf{SL}(2, \mathbf{R}) / \mathrm{N}\) 与 \(\mathbf{R}^2 - \{0\}\) 通过映射 \(g \mapsto ge_1\) 等同，其中 \(e_1\) 是 \(\mathbf{R}^2\) 的典范基的第一个向量；齐性空间 \(\mathbf{SL}(2, \mathbf{R}) / \mathrm{P}\) 与子空间 \((\mathbf{R} - \{0\}) \times \{0\}\) （\(\mathbf{SL}(2, \mathbf{R}) / \mathrm{N}\) 的）等同。
\item
  G 上的每个连续函数 \(f\) ，若满足 \(f(ngn') = f(g)\) （对所有 \((n, n') \in \mathrm{N}^2\) ），则满足 \(f(pgp') = f(g)\) （对所有 \((p, p') \in \mathrm{P}^2\) ）。（把 \(f\) 解释为连续函数 \(\varphi\) （在 \(\mathbf{R}^2 - \{0\}\) 上、在 N 的自然作用下不变），并验证它在 \(\mathbf{SL}(2, \mathbf{R}) / \mathrm{P}\) 上是常数。）
\end{enumerate}

\begin{enumerate}
\def\labelenumi{\alph{enumi})}
\setcounter{enumi}{2}
\tightlist
\item
  有 \(\pi(g)x = x\) 对所有 \(g \in G\) 成立。
\end{enumerate}

\begin{enumerate}
\def\labelenumi{\arabic{enumi})}
\setcounter{enumi}{18}
\tightlist
\item
  设 \(n \geqslant 2\) 为整数。考虑实李群 \(\mathrm{G} = \mathbf{SL}(n, \mathbf{R})\) 。用 N 与 H 表示 G 中元素分别为分块形式
\end{enumerate}

\[
\left( \begin{array}{c c} \operatorname{Id} _ {n - 1} & x \\ 0 & 1 \end{array} \right), \quad (x \in \mathbf {R} ^ {n - 1}), \qquad \left( \begin{array}{c c} g & 0 \\ 0 & 1 \end{array} \right), \quad (g \in \mathbf {S L} (n - 1, \mathbf {R}))
\]

的闭子群。群 N 与 \(R^{n-1}\) 等同，群 H 与 \(\mathbf{SL}(n-1,\mathbf{R})\) 等同。

\begin{enumerate}
\def\labelenumi{\alph{enumi})}
\item
  设 \(\pi\) 为 G 在复希尔伯特空间 E 中的酉表示。E 的元素 \(x\) 在 G 下不变，当且仅当它在 N 下不变。
\item
  若 \(n \geqslant 3\) ，则群 G 具有性质 ]T[。（把上一问题的结果与上一习题结合起来。）
\item
  若 \(n \geqslant 3\) ，则群 \(\widetilde{G} = R^{n} \rtimes G\) 具有性质 ]T[。（设 \(\pi\) 为满足 \(1_{\widetilde{G}} \leqslant \pi\) 的酉表示；由上一问题，存在非零向量 \(x \neq 0\) ，它被 \(\widetilde{G}\) 的子群 G 保持不变；证明相应的矩阵系数是常数。）
\item
  若 \(n \geqslant 3\) ，则群 \(\mathbf{SL}(n, \mathbf{Z})\) 具有性质 ]T[。（利用 INT, VII, p.~116 的习题 7。）
\item
  设 \(n \geqslant 3\) 。群 \(\mathbf{SL}(n, \mathbf{Z})\) 是有限生成的。设 H 为 \(\mathbf{SL}(n, \mathbf{Z})\) 的有限指标正规子群所成的族；它是无限的。对 \(\mathbf{SL}(n, \mathbf{Z})\) 的每个有限生成子集 S，凯莱图族 \((\mathcal{G}(\mathbf{SL}(n, \mathbf{Z})/\mathrm{H}, \mathrm{SH}/\mathrm{H}))_{\mathrm{H} \in \mathcal{H}}\) 是一个扩展图族（参见 TA, 2, p.~221, 习题 7 及 IV, p.~323 的习题 16）。
\end{enumerate}

\begin{enumerate}
\def\labelenumi{\arabic{enumi})}
\setcounter{enumi}{19}
\tightlist
\item
  设 G 为局部紧群，\(\mu\) 为 G 上的左哈尔测度。我们用 \(\Delta\) 表示 G 的模。对 G 上的任意函数 f，我们将记 \(\widetilde{f} = \Delta f^{*}\) 。
\end{enumerate}

G 上的测度 \(\nu\) 称为正定型的，如果 \(\langle \nu, \tilde{f} * f \rangle \geqslant 0\) 对每个函数 \(f \in \mathcal{K}(\mathrm{G})\) 成立。G 上的局部可积函数 \(\varphi\) 称为正定型的，如果 \(\langle \varphi, f^* * f \rangle \geqslant 0\) 对每个函数 \(f \in \mathcal{K}(\mathrm{G})\) 成立。

\begin{enumerate}
\def\labelenumi{\alph{enumi})}
\item
  假设 \(\nu\) 有界。测度 \(\nu\) 是正定型的，当且仅当 \(\gamma(\nu) \geqslant 0\) 。
\item
  设 \(\varphi\in\mathcal{K}(\mathrm{G})\) 。函数 \(\varphi\) 作为局部可积函数是正定型的，当且仅当它按照 V, p.~442 的定义 4 是正定型的。
\item
  设 \(\nu\) 为 G 上的一个正定型测度。公式
\end{enumerate}

\[
(f \mid g) _ {\nu} = \langle \nu , \widetilde {f} * g \rangle
\]

在 \(\mathcal{K}(\mathrm{G})\) 上定义了一个正埃尔米特型。我们用 \(\mathrm{E}_{\nu}\) 表示赋予了该埃尔米特型的 \(\mathcal{K}(\mathrm{G})\) 的分离完备希尔伯特空间。

\begin{enumerate}
\def\labelenumi{\alph{enumi})}
\setcounter{enumi}{3}
\tightlist
\item
  对 \(f \in \mathcal{K}(\mathrm{G})\) ，函数 \(\sigma(g)f \colon g \mapsto (\varepsilon_g * f)\Delta^{1/2}\) 属于 \(\mathcal{K}(\mathrm{G})\) ，且 \(\| \sigma(g)f \|_{\nu} = \|f\|_{\nu}\) 。映射 \(g \mapsto \sigma(g)\) 延拓为 G 在 \(\mathrm{E}_{\nu}\) 中的一个酉表示。
\end{enumerate}

\begin{enumerate}
\def\labelenumi{\alph{enumi})}
\setcounter{enumi}{4}
\tightlist
\item
  设 U 为 G 中 e 的开邻域，\(\nu\) 为 G 上的正定型测度，使得存在实数 \(c \geqslant 0\) 满足
\end{enumerate}

\[
\langle \nu , f * \widetilde {f} \rangle \leqslant c \left(\int_ {\mathrm{G}} f d \mu\right) ^ {2}
\]

  其中 \(f \in \mathcal{K}(\mathrm{G})\) 为正的且在 U 外为零。存在 G 上的连续正定型函数 \(\varphi\) ，使得 \(\nu = \varphi \Delta^{-1/2} \cdot \mu\) 。（设 \(j\) 为从 \(\mathcal{K}(\mathrm{G})\) 到 \(\mathrm{E}_{\nu}\) 的典范态射；若 \(\mathfrak{B}\) 是 \(e\) 的邻域滤子的一组基（其元素都含于 U），而 \(f_{\mathrm{V}} \in \mathcal{K}_{+}(\mathrm{G})\) 是支集含于 \(V \in B\) 、积分为 1 的函数，证明 \(j(f_{\mathrm{V}})\) 沿 B 狭收敛于元素 \(y_{0}\) （属于 \(E_{\nu}\) ），使得

\[
(j (f) \mid y _ {0}) = \int_ {\mathrm{G}} f ^ {*} d \mu
\]

对 \(f \in \mathcal{K}(\mathrm{G})\) 成立；令 \(\varphi(g) = (y_0 \mid \sigma(g)y_0)\) ，验证 \(\varphi\) 具有所需的性质。

\begin{enumerate}
\def\labelenumi{\alph{enumi})}
\setcounter{enumi}{5}
\tightlist
\item
  设 \(\varphi_{1}\) 为 G 上的连续正定函数，\(\varphi_{2}\) 为局部可积的正定型函数。若 \(\varphi_{1}-\varphi_{2}\) 是正定型的，则 \(\varphi_{2}\) 局部几乎处处等于一个连续的正定型函数。（验证可以应用上一问题。）
\end{enumerate}

\begin{enumerate}
\def\labelenumi{\arabic{enumi})}
\setcounter{enumi}{20}
\tightlist
\item
  我们沿用上一习题的记号。
\end{enumerate}

\begin{enumerate}
\def\labelenumi{\alph{enumi})}
\item
  设 \(f \in \mathcal{L}^{2}(G)\) 作为局部可积函数是正定型的。\(L^{2}(G)\) 上以 \(\mathcal{K}(G)\) 为定义域、由 \(\varphi \mapsto \gamma(\varphi)f\) 定义的部分算子是正的。我们用 \(\varrho(f)\) 表示 f 的弗里德里希斯延拓（IV, p.~295 的例 8）；它是 \(L^{2}(G)\) 上的一个正自伴部分算子。
\item
  设 \(f \in \mathcal{L}^2(\mathrm{G})\) 。我们称 \(f\) 是清醒的，如果存在实数 \(c \geqslant 0\) ，使得 \(\| \boldsymbol{\gamma}(\varphi)f \| \leqslant c \| \varphi \|\) 对所有 \(\varphi \in \mathcal{K}(\mathrm{G}) \subset \mathcal{L}^2(\mathrm{G})\) 成立。这时我们用 \(\varrho(f)\) 表示 \(\mathrm{L}^2(\mathrm{G})\) 中延拓 \(\varphi \mapsto \boldsymbol{\gamma}(\varphi)f\) 的自同态。若 \(f\) 是正定型的，这一定义与上一个定义一致，并且这一条件此时等价于说 \(\varrho(f)\) 是正的。
\item
  设 \(f_{1}\) 与 \(f_{2}\) 为 \(\mathcal{L}^{2}(G)\) 中的清醒函数。若 \(\varrho(f_{1}) = \varrho(f_{2})\) ，则 \(f_{1}\) 与 \(f_{2}\) 在 \(L^{2}(G)\) 中相等。
\end{enumerate}

\begin{enumerate}
\def\labelenumi{\alph{enumi})}
\setcounter{enumi}{3}
\tightlist
\item
  设 \(f \in \mathcal{L}^2(\mathrm{G})\) 为正定型的。对所有 \(g \in \mathrm{G}\) ，我们有 \(\varrho(f) \circ \boldsymbol{\gamma}(g) = \boldsymbol{\gamma}(g) \circ \varrho(f)\) 。此外，对每个 \(h \in \operatorname{dom}(\varrho(f))\) ，我们有 \(\varrho(f)h = h * f\) 。
\end{enumerate}

\begin{enumerate}
\def\labelenumi{\alph{enumi})}
\setcounter{enumi}{4}
\tightlist
\item
  设 \(f_{1}\) 与 \(f_{2}\) 为 \(\mathcal{L}^{2}(\mathrm{G})\) 中的清醒元素，使得 \(\varrho(f_{1})\) 与 \(\varrho(f_{2})\) 交换。函数 \(f_{1}*f_{2}\) 是 G 上的连续正定函数；它是清醒的，并且我们有
\end{enumerate}

\[
\varrho (f _ {1} * f _ {2}) = \varrho (f _ {1}) \varrho (f _ {2}), \quad \langle f _ {1} | f _ {2} \rangle \geqslant 0.
\]

\begin{enumerate}
\def\labelenumi{\alph{enumi})}
\setcounter{enumi}{5}
\tightlist
\item
  在上一问题的记号与假设下，若 \(f_{2} - f_{1}\) 是正定型的，则
\end{enumerate}

\[
\left\| f _ {2} - f _ {1} \right\| ^ {2} \leqslant \left\| f _ {2} \right\| ^ {2} - \left\| f _ {1} \right\| ^ {2}.
\]

\begin{enumerate}
\def\labelenumi{\alph{enumi})}
\setcounter{enumi}{6}
\tightlist
\item
  设 \((f_n)_{n\in \mathbf{N}}\) 为 \(\mathcal{L}^2 (\mathrm{G})\) 中两两交换的清醒正定型元素构成的序列，且 \(f_{n + 1} - f_n\) 对所有 \(n\in \mathbf{N}\) 是正定型的。若序列 \((\| f\| _n)\) 有界，则序列 \((f_n)\) 在 \(\mathcal{L}^2 (\mathrm{G})\) 中收敛。
\end{enumerate}

\begin{enumerate}
\def\labelenumi{\arabic{enumi})}
\setcounter{enumi}{21}
\tightlist
\item
  我们沿用上一习题的记号。
\end{enumerate}

设 \(f\) 为属于 \(\mathcal{L}^2 (\mathrm{G})\) 的连续正定型函数。

\begin{enumerate}
\def\labelenumi{\alph{enumi})}
\item
  存在 [0,1] 上多项式函数的递增序列 \((p_{n})_{n\in\mathbf{N}}\) ，它在 [0,1] 上一致收敛于函数 \(t\mapsto\sqrt{t}\) 。
\item
  假设 \(f\) 是清醒的，且 \(0 \leqslant f \leqslant 1\) （在 \(\mathcal{L}(\mathrm{L}^2(\mathrm{G}))\) 中）。存在序列 \((h_n)\) （在 \(\mathcal{L}^2(\mathrm{G})\) 中），满足下列条件：
\end{enumerate}

\begin{enumerate}
\def\labelenumi{(\roman{enumi})}
\item
  每个 \(h_n\) 都是清醒的且是正定型的；
\item
  函数 \(h_{n+1} - h_n\) 对所有 \(n \in \mathbf{N}\) 是正定型的；
\item
  自同态 \(\varrho(h_n)\) 两两交换；
\item
  有 \(\varrho(h_n)^2 \leqslant \varrho(f)\) （对所有 \(n\) ）。
\end{enumerate}

\begin{enumerate}
\def\labelenumi{\alph{enumi})}
\setcounter{enumi}{2}
\tightlist
\item
  序列 \((h_n)_{n\in \mathbf{N}}\) 在 \(\mathcal{L}^2 (\mathrm{G})\) 中收敛；若 \(h\) 是 \((h_n)\) 的一个极限，则 \(h\) 是正定型且清醒的，并满足 \(\varrho (h)^2 = \varrho (f)\) ；我们有 \(f = h*h\) 。
\end{enumerate}

\begin{enumerate}
\def\labelenumi{\alph{enumi})}
\setcounter{enumi}{3}
\tightlist
\item
  我们不再假设 \(f\) 是清醒的。对 \(t \in \mathbf{R}\) ，用 \(p_t\) 表示部分自伴算子 \(\varrho(f)\) 的由 \([- \infty, t]\) 定义的谱投影。我们有 \(p_t \circ \boldsymbol{\gamma}(g) = \boldsymbol{\gamma}(g) \circ p_t\) 对所有 \(t \in \mathbf{R}\) 与 \(g \in G\) 成立。
\end{enumerate}

\begin{enumerate}
\def\labelenumi{\alph{enumi})}
\setcounter{enumi}{4}
\tightlist
\item
  对 \(t \in \mathbf{R}\) ，令 \(f_t = p_t(f)\) ；它是 \(\mathrm{L}^2(\mathrm{G})\) 中某个连续正定的清醒函数（属于 \(\mathcal{L}^2(\mathrm{G})\) ）的类。（验证 \(\gamma(\varphi)f_t = \varrho(f)p_t(g)\) 对 \(g \in \mathcal{H}(\mathrm{G})\) 成立，由此推出 \(\varrho(f_t) = \varrho(f) \circ p_t\) ，再应用习题 20 的结果。）
\end{enumerate}

\begin{enumerate}
\def\labelenumi{\alph{enumi})}
\setcounter{enumi}{5}
\tightlist
\item
  存在 \(h \in \mathcal{L}^2(\mathrm{G})\) ，使得 \(f = h * h\) 。（考虑 \(f_n = p_n(f)\) （\(n \in \mathbf{N}\) ），并应用部分 \(c\) 。）
\end{enumerate}

\begin{enumerate}
\def\labelenumi{\arabic{enumi})}
\setcounter{enumi}{22}
\tightlist
\item
  我们记 \(G = R_{+}^{*}\) ；它是一个局部紧交换群，且测度 \(\mu = dx/x\) 是 G 上的哈尔测度。群 G 与群 iR 通过映射 \((x, it) \mapsto x^{it}\) 处于对偶；\(\mu\) 的对偶测度是 iR 上的勒贝格测度（参见 II, p. 236 的推论 4）。我们用 F 表示 G 的傅里叶变换，用 \(\overline{F}\) 表示 \(\widehat{G}\) 的傅里叶余变换。
\end{enumerate}

我们用 X 表示 G 上的恒等函数。

\begin{enumerate}
\def\labelenumi{\alph{enumi})}
\item
  设 \(\nu\) 为 G 上的正测度。我们用 \(\mathrm{I}_{\nu}\) 表示这样的 \(\sigma \in \mathbf{R}\) 所成的集合：函数 \(\mathrm{X}^{\sigma}\) 是 \(\nu\) -可积的；用 \(\mathrm{B}_{\nu}\) 表示这样的 \(s \in \mathbf{R}\) 所成的集合：\(\mathcal{R}(s) \in \mathrm{I}_{\nu}\) 。集合 \(\mathrm{I}_{\nu}\) 是 \(\mathbf{R}\) 的一个区间，而函数 \(s \mapsto \nu(\mathrm{X}^s)\) 在 \(\mathrm{B}_{\nu}\) 的内部有定义且全纯。
\item
  设 \(f\) 为 G 上的局部可积函数。我们记 \(\mathrm{I}_{f} = \mathrm{I}_{f\cdot \mu}\) 与 \(\mathrm{B}_f = \mathrm{B}_{f\cdot \mu}\) 。对每个 \(s \in \mathrm{B}_f\) ，我们令
\end{enumerate}

\[
\widehat {f} (s) = \int_ {\mathrm{G}} f (x) x ^ {s} d \mu (x).
\]

若存在 \(\delta > 0\) 与 \(c \geqslant 0\) ，使得 \(|\widehat{f}(s)| \leqslant c(1 + |s|)^{-1 - \delta}\) 对 \(s \in \mathrm{B}_f\) 成立，则对所有 \(\sigma\) （属于 \(\mathrm{B}_f\) 的内部），我们有

\[
f (x) = \frac {1}{2 \pi} \int_ {\mathbf {R}} \widehat {f} (\sigma + i t) x ^ {- \sigma - i t} d t.
\]

\begin{enumerate}
\def\labelenumi{\alph{enumi})}
\setcounter{enumi}{2}
\tightlist
\item
  设 \(f \in \mathcal{L}^{2}(\mathrm{G})\) 。假设 0 属于 \(I_{f}\) 的内部。则我们有 \(\overline{\mathcal{F}} f(it) = \widehat{f}(it)\) 对几乎所有 \(t \in \mathbf{R}\) 成立。
\end{enumerate}

\begin{enumerate}
\def\labelenumi{\alph{enumi})}
\setcounter{enumi}{3}
\item
  设 \(f\) 与 \(g\) 为 \(\mathbf{G}\) 上的可测函数。设 \(\sigma\) 为实数，使得 \(\mathrm{X}^{-2\sigma}f\in \mathrm{L}^2 (\mathrm{G})\) 且 \(\mathrm{X}^{2\sigma}g\in \mathrm{L}^2 (\mathrm{G})\) 。则我们有

\[
\int_ {\mathrm{G}} f g d \mu = \frac {1}{2 \pi} \int_ {\mathbf {R}} \widehat {f} (- \sigma - i t) \widehat {g} (\sigma + i t) d t.
\]

\end{enumerate}

\begin{enumerate}
\def\labelenumi{\alph{enumi})}
\setcounter{enumi}{4}
\tightlist
\item
  设 \(f\) 与 \(g\) 为 \(G\) 上的可测函数。假设
\end{enumerate}

\begin{enumerate}
\def\labelenumi{(\roman{enumi})}
\item
  存在 \(c \in \mathbf{R}\) ，使得 \(|f(x)| \leqslant cx / (1 + x)\) 对所有 \(x \in \mathbf{G}\) 成立；
\item
  有 ]0,1[ ⊂ I \(_{g}\) ，且对每个满足 0 < σ < 1 的 σ，函数 X \(^{\sigma}\) g 有界；
\item
  存在 \(c' \in \mathbf{R}\) ，使得 \(|\widehat{g}(s)| \leqslant c'(1 + |s|)^{-2}\) 对所有 \(s \in \mathbf{C}\) （满足 \(\mathcal{R}(s) \in ]0, 1[\) ）成立。
\end{enumerate}

则我们有

\[
\int_ {\mathrm{G}} f g d \mu = \frac {1}{2 \pi} \int_ {\mathbf {R}} \widehat {f} (- \sigma - i t) \widehat {g} (\sigma + i t) d t
\]

对所有 \(\sigma\in\mathbf{R}\) （满足 \(0<\sigma<1\) ）成立。

\begin{enumerate}
\def\labelenumi{\alph{enumi})}
\setcounter{enumi}{5}
\tightlist
\item
  设 \(\sigma \in \mathbf{R}\) ，\(f \in \mathcal{C}(\sigma + i\mathbf{R})\) 。G 上存在正测度 \(\nu\) ，使得 \(\sigma \in \mathrm{I}_{\nu}\) 且 \(\widehat{\nu}(\sigma + it) = f(\sigma + it)\) 对所有 \(t \in \mathbf{R}\) 成立，当且仅当对每个整数 \(n \geqslant 1\) 、复数的每个族 \((t_i)_{0 \leqslant i \leqslant n}\) 与任意族 \((z_i)_{0 \leqslant i \leqslant n}\) （\(\mathcal{R}(z_i) = \sigma\) 对所有 \(i\) ），我们有
\end{enumerate}

\[
\sum_ {i = 0} ^ {n} \sum_ {j = 0} ^ {n} \bar {t} _ {i} t _ {j} f \left(\frac {1}{2} (\overline {{z}} _ {i} + z _ {j})\right) \geqslant 0.
\]

